\documentclass[a4paper,11pt]{article}
        \usepackage[top=15mm,bottom=18mm,left=16mm,right=16mm]{geometry}
        \usepackage{fontspec}
        \usepackage{xeCJK}
        \setmainfont{Times New Roman}
        \setCJKmainfont{Yu Gothic}
        \setmonofont{Consolas}
        \setCJKmonofont{Yu Gothic}
        \usepackage{booktabs}
        \usepackage{longtable}
        \usepackage{tabularx}
        \usepackage{array}
        \usepackage{enumitem}
        \usepackage{hyperref}
        \usepackage{listings}
        \usepackage{xcolor}
        \hypersetup{
          colorlinks=true,
          linkcolor=blue,
          urlcolor=blue,
          pdftitle={GPU 候補の厳密検証},
          pdfauthor={Codex}
        }
        \lstset{
          basicstyle=\ttfamily\small,
          breaklines=true,
          columns=fullflexible,
          keepspaces=true,
          frame=single,
          framerule=0.2pt,
          rulecolor=\color{black},
          showstringspaces=false
        }
        \setlength{\parskip}{0.42em}
        \setlength{\parindent}{1em}
        \renewcommand{\arraystretch}{1.15}
        \title{39. GPU 候補の厳密検証}
        \author{solar\_ws0129-main}
        \date{2026-07-29}
        \begin{document}
        \maketitle

        \section*{対象}
        \begin{tabularx}{\linewidth}{>{\raggedright\arraybackslash}p{0.18\linewidth} X}
        \toprule
        項目 & 内容 \\
        \midrule
        ファイル & \path|scripts/validate_gpu_upper_policy_candidates.py| \\
        種別 & Python \\
        区分 & planning validation \\
        役割 & GPU 探索で得た speed policy 候補を CPU 側 exact replay と gate で再判定する。 \\
        起動文脈 & GPU acceptance の中心。 \\
        \bottomrule
        \end{tabularx}

        \section*{このファイルを読む理由}
        GPU 探索で得た speed policy 候補を CPU 側 exact replay と gate で再判定する。

        \begin{itemize}[leftmargin=1.6em]
        \item numerical match、mission feasibility、gate pass を確認する。
        \end{itemize}

        \section*{呼び出し関係}
        \subsection*{どこから呼ばれるか}
        \begin{itemize}[leftmargin=1.6em]
        \item GPU acceptance pipeline
\item 手動検証
        \end{itemize}

        \subsection*{次に読むと理解が進むファイル}
        \begin{itemize}[leftmargin=1.6em]
        \item \path|scripts/run_upper_mesh_convergence.py|
\item \path|scripts/solar_sim.py|
        \end{itemize}

        \subsection*{同一リポジトリ内の主要依存}
        \begin{itemize}[leftmargin=1.6em]
        \item \path|scripts/run_upper_mesh_convergence.py|
        \end{itemize}

        \section*{ソース冒頭の把握}
        モジュール docstring または先頭意図の要約:

        Replay every GPU policy in the exact 1 Hz simulator and select the fastest feasible one.

        \subsection*{代表コード断片}
        \begin{lstlisting}[language=Python]
ROOT = Path(__file__).resolve().parents[1]


def resolve_result_path(raw: str) -> Path:
    path = Path(str(raw))
    return path if path.is_absolute() else (ROOT / path).resolve()


def resolve_profile_asset(profile_path: Path, raw: str) -> Path:
    path = Path(str(raw))
    return path if path.is_absolute() else (profile_path.parent / path).resolve()


def evaluate_event_timing(detail: pd.DataFrame, cfg: dict, profile_path: Path) -> dict:
    """Check official control-stop closing times and the absolute finish deadline."""
    paths = cfg.get("paths", {}) or {}
    simulation = cfg.get("simulation", {}) or {}
    stop_path = resolve_profile_asset(profile_path, paths.get("stop_yaml", ""))
    stop_cfg = (
        yaml.safe_load(stop_path.read_text(encoding="utf-8-sig")) or {}
        if stop_path.is_file()
        else {}
\end{lstlisting}


        \section*{主要構造の読み取り}
        主要関数は resolve\_result\_path, resolve\_profile\_asset, evaluate\_event\_timing, exact\_replay\_signature, inspect\_prediction\_mesh, prediction\_execution\_soc\_errors, evaluate\_soc\_guard\_intervention, main。 CLI 引数宣言は 8 件。

        \subsection*{主要クラス・関数}
            \begin{longtable}{p{0.07\linewidth} p{0.16\linewidth} p{0.25\linewidth} p{0.40\linewidth}}
            \toprule
            行 & 種別 & 名前 & 補足 \\
            \midrule
            26 & function & \path|resolve_result_path| & args: raw \\
31 & function & \path|resolve_profile_asset| & args: profile\_path, raw \\
36 & function & \path|evaluate_event_timing| & args: detail, cfg, profile\_path \\
112 & function & \path|exact_replay_signature| & args: profile, policy \\
137 & function & \path|inspect_prediction_mesh| & args: manifest, manifest\_path \\
199 & function & \path|prediction_execution_soc_errors| & args: manifest \\
228 & function & \path|evaluate_soc_guard_intervention| & args: detail, manifest \\
249 & function & \path|main| &  \\
            \bottomrule
            \end{longtable}

        \subsection*{ファイルを上から読んだときの定義順}
            \begin{longtable}{p{0.07\linewidth} p{0.84\linewidth}}
            \toprule
            行 & 内容 \\
            \midrule
            17 & 例外処理を伴う try ブロックを実行する。 \\
23 & ROOT に Path(\_\_file\_\_).resolve().parents[1] の結果を代入する。 \\
26 & 関数 resolve\_result\_path を定義する。 \\
31 & 関数 resolve\_profile\_asset を定義する。 \\
36 & 関数 evaluate\_event\_timing を定義する。 \\
112 & 関数 exact\_replay\_signature を定義する。 \\
137 & 関数 inspect\_prediction\_mesh を定義する。 \\
199 & 関数 prediction\_execution\_soc\_errors を定義する。 \\
228 & 関数 evaluate\_soc\_guard\_intervention を定義する。 \\
249 & 関数 main を定義する。 \\
476 & 条件 \_\_name\_\_ == '\_\_main\_\_' を判定し、真なら内部処理を行う。 \\
            \bottomrule
            \end{longtable}

        \subsection*{import 群の詳細}
            \begin{longtable}{p{0.07\linewidth} p{0.33\linewidth} p{0.50\linewidth}}
            \toprule
            行 & import 文 & このファイルで読む理由 \\
            \midrule
            4 & \path|from __future__ import annotations| & 型ヒントの遅延評価を有効にし、自己参照型や forward reference を安全に扱うため。 このファイル内での使用位置は少ないか、間接利用である。 \\
6 & \path|import argparse| & CLI 引数を宣言し、実行時パラメータを外から受け取るため。 このファイル内での主な使用位置は L250。 \\
7 & \path|import json| & manifest、checkpoint、UDP payload をやり取りするため。 このファイル内での主な使用位置は L309, L340, L470, L472。 \\
8 & \path|import math| & clip、sqrt、sin/cos、有限判定など数値ロジックに使うため。 このファイル内での主な使用位置は L148, L205, L210。 \\
9 & \path|import shutil| & shutil モジュールを利用するため。 このファイル内での主な使用位置は L458。 \\
10 & \path|import subprocess| & git 情報取得や外部コマンド実行を行うため。 このファイル内での主な使用位置は L318, L322。 \\
11 & \path|import sys| & sys モジュールを利用するため。 このファイル内での主な使用位置は L319。 \\
12 & \path|from pathlib import Path| & ファイルやディレクトリを安全に扱うため。 このファイル内での主な使用位置は L23, L26, L27, L31, L32, L36, L112, L123, ...。 \\
14 & \path|import pandas as pd| & CSV や時刻列の読込・整形・再サンプリングに使うため。 このファイル内での主な使用位置は L36, L48, L49, L50, L63, L64, L88, L89, ...。 \\
15 & \path|import yaml| & profile、stop、schedule、summary YAML を読み書きするため。 このファイル内での主な使用位置は L42, L127, L272, L297。 \\
18 & \path|from scripts.run_upper_mesh_convergence import file_sha256, materialize_profile| & upper mesh 収束確認 から file\_sha256, materialize\_profile を読み込み、このファイルの内部処理を分担させるため。 実体ファイルは scripts/run\_upper\_mesh\_convergence.py。 このファイル内での主な使用位置は L125, L126, L133, L282。 \\
20 & \path|from run_upper_mesh_convergence import file_sha256, materialize_profile| & run\_upper\_mesh\_convergence から file\_sha256, materialize\_profile を読み込み、このファイルの処理を組み立てるため。 このファイル内での主な使用位置は L125, L126, L133, L282。 \\
113 & \path|import hashlib| & snapshot ID や入力資産の digest を作るため。 このファイル内での主な使用位置は L134。 \\
            \bottomrule
            \end{longtable}

        \section*{関数・クラスを上から順に解説}
        \subsection*{L26 関数 \path|resolve_result_path|}
            \begin{itemize}[leftmargin=1.6em]
              \item 定義: \path|resolve_result_path(raw)|
              \item docstring: 明示されていない。
              \item このブロックが直接呼ぶ主な関数・メソッド: Path, is\_absolute, resolve, str
              \item 戻り値の要点: path if path.is\_absolute() else (ROOT / path).resolve()
            \end{itemize}
            \begin{enumerate}[leftmargin=1.8em]
            \item path に Path(str(raw)) の結果を代入する。
\item path if path.is\_absolute() else (ROOT / path).resolve() を返す。
            \end{enumerate}

\subsection*{L31 関数 \path|resolve_profile_asset|}
            \begin{itemize}[leftmargin=1.6em]
              \item 定義: \path|resolve_profile_asset(profile_path, raw)|
              \item docstring: 明示されていない。
              \item このブロックが直接呼ぶ主な関数・メソッド: Path, is\_absolute, resolve, str
              \item 戻り値の要点: path if path.is\_absolute() else (profile\_path.parent / path).resolve()
            \end{itemize}
            \begin{enumerate}[leftmargin=1.8em]
            \item path に Path(str(raw)) の結果を代入する。
\item path if path.is\_absolute() else (profile\_path.parent / path).resolve() を返す。
            \end{enumerate}

\subsection*{L36 関数 \path|evaluate_event_timing|}
            \begin{itemize}[leftmargin=1.6em]
              \item 定義: \path|evaluate_event_timing(detail, cfg, profile_path)|
              \item docstring: Check official control-stop closing times and the absolute finish deadline.
              \item このブロックが直接呼ぶ主な関数・メソッド: DataFrame, Timestamp, all, append, bool, dropna, float, get, is\_file, isna, isoformat, max
              \item 戻り値の要点: \{'control\_stop\_windows\_pass': bool(all\_stops\_observed and all((row['passed'] for row in stop\_results))), 'finish\_deadline\_pass': bool(deadline\_late\_sec <= 1.0), 'max\_control\_stop\_late\_sec': max\_late\_sec, 'finish\_deadline\_late\_sec': deadline\_late\_sec, 'finish\_time\_utc': finish\_time.isoformat() if not pd.isna(finish\_time) else '', 'race\_deadline\_utc': deadline.isoformat() if deadline\_raw else '', 'control\_stops': stop\_results\}
            \end{itemize}
            \begin{enumerate}[leftmargin=1.8em]
            \item paths に cfg.get('paths', \{\}) or \{\} の結果を代入する。
\item simulation に cfg.get('simulation', \{\}) or \{\} の結果を代入する。
\item stop\_path に resolve\_profile\_asset(profile\_path, paths.get('stop\_yaml', '')) の結果を代入する。
\item stop\_cfg に yaml.safe\_load(stop\_path.read\_text(encoding='utf-8-sig')) or \{\} if stop\_path.is\_file() else \{\} の結果を代入する。
\item distance\_col に 's\_end\_km' if 's\_end\_km' in detail else 's\_km' の結果を代入する。
\item time\_col に 'time\_end\_utc' if 'time\_end\_utc' in detail else 'time\_utc' の結果を代入する。
\item distance に pd.to\_numeric(detail.get(distance\_col), errors='coerce') の結果を代入する。
\item times に pd.to\_datetime(detail.get(time\_col), errors='coerce', utc=True, format='mixed') の結果を代入する。
\item ordered に pd.DataFrame(\{'s\_km': distance, 'time': times\}).dropna().sort\_values('time') の結果を代入する。
\item stop\_results に [] の結果を代入する。
            \end{enumerate}

\subsection*{L112 関数 \path|exact_replay_signature|}
            \begin{itemize}[leftmargin=1.6em]
              \item 定義: \path|exact_replay_signature(profile, policy)|
              \item docstring: 明示されていない。
              \item このブロックが直接呼ぶ主な関数・メソッド: Path, encode, file\_sha256, get, hexdigest, is\_file, isinstance, join, read\_text, resolve, resolve\_profile\_asset, safe\_load
              \item 戻り値の要点: hashlib.sha256(payload.encode('ascii')).hexdigest()
            \end{itemize}
            \begin{enumerate}[leftmargin=1.8em]
            \item Import 文を実行する。
\item dependencies に (ROOT / 'scripts' / 'solar\_sim.py', ROOT / 'mpc\_solarcar' / 'upper\_horizon.py', ROOT / 'mpc\_solarcar' / 'upper\_solver.py', ROOT / 'mpc\_solarcar' / 'upper\_policy.py', ROOT / 'mpc\_solarcar' / 'upper\_cost.py', ROOT / 'mpc\_solarcar' / 'model.py', ROOT / 'mpc\_solarcar' / 'signal\_utils.py', Path(\_\_file\_\_).resolve()) の結果を代入する。
\item payload に file\_sha256(profile) + file\_sha256(policy) の結果を代入する。
\item payload を Add で更新する。
\item cfg に yaml.safe\_load(profile.read\_text(encoding='utf-8-sig')) or \{\} の結果を代入する。
\item (cfg.get('paths', \{\}) or \{\}).values() を順に走査し、各要素を raw に入れて処理する。
\item hashlib.sha256(payload.encode('ascii')).hexdigest() を返す。
            \end{enumerate}

\subsection*{L137 関数 \path|inspect_prediction_mesh|}
            \begin{itemize}[leftmargin=1.6em]
              \item 定義: \path|inspect_prediction_mesh(manifest, manifest_path, requested_ds_km, start_s_km, race_km)|
              \item docstring: Verify that the selected prediction trace used the requested distance mesh.
              \item このブロックが直接呼ぶ主な関数・メソッド: Path, abs, bool, ceil, dropna, float, get, int, is\_file, len, list, max
              \item 戻り値の要点: result / result / result / result
            \end{itemize}
            \begin{enumerate}[leftmargin=1.8em]
            \item requested\_ds\_km に float(requested\_ds\_km) の結果を代入する。
\item remaining\_km に max(0.0, float(race\_km) - float(start\_s\_km)) の結果を代入する。
\item expected\_min\_steps に int(math.ceil(remaining\_km / requested\_ds\_km - 1e-09)) の結果を代入する。
\item result に \{'valid': False, 'expected\_min\_steps': expected\_min\_steps, 'prediction\_steps': 0, 'trace\_drive\_steps': 0, 'max\_trace\_ds\_km': float('nan'), 'trace\_terminal\_km': float('nan')\} の結果を代入する。
\item diagnostics に list(manifest.get('upper\_solver\_diagnostics', []) or []) の結果を代入する。
\item 条件 not diagnostics を判定し、真なら内部処理を行う。
\item diagnostic に diagnostics[0] の結果を代入する。
\item result['prediction\_steps'] に int(diagnostic.get('prediction\_steps', 0) or 0) の結果を代入する。
\item raw\_trace に str(diagnostic.get('selected\_prediction\_trace\_csv', '') or '') の結果を代入する。
\item 条件 not raw\_trace を判定し、真なら内部処理を行う。
            \end{enumerate}

\subsection*{L199 関数 \path|prediction_execution_soc_errors|}
            \begin{itemize}[leftmargin=1.6em]
              \item 定義: \path|prediction_execution_soc_errors(manifest)|
              \item docstring: 明示されていない。
              \item このブロックが直接呼ぶ主な関数・メソッド: append, dict, float, get, int, isfinite, list
              \item 戻り値の要点: \{'initial\_prediction\_soc': initial\_prediction, 'latest\_nontrivial\_prediction\_soc': latest\_prediction, 'initial\_error': final\_soc - initial\_prediction, 'latest\_nontrivial\_error': final\_soc - latest\_prediction\} / \{'initial\_prediction\_soc': float('nan'), 'latest\_nontrivial\_prediction\_soc': float('nan'), 'initial\_error': float('inf'), 'latest\_nontrivial\_error': float('inf')\}
            \end{itemize}
            \begin{enumerate}[leftmargin=1.8em]
            \item diagnostics に list(manifest.get('upper\_solver\_diagnostics', []) or []) の結果を代入する。
\item predictions に [] の結果を代入する。
\item diagnostics を順に走査し、各要素を diagnostic に入れて処理する。
\item final\_soc に float(manifest.get('final\_soc', float('nan'))) の結果を代入する。
\item 条件 not predictions or not math.isfinite(final\_soc) を判定し、真なら内部処理を行う。
\item initial\_prediction に predictions[0][1] の結果を代入する。
\item nontrivial に [soc for steps, soc in predictions if steps > 1] の結果を代入する。
\item latest\_prediction に nontrivial[-1] if nontrivial else predictions[-1][1] の結果を代入する。
\item \{'initial\_prediction\_soc': initial\_prediction, 'latest\_nontrivial\_prediction\_soc': latest\_prediction, 'initial\_error': final\_soc - initial\_prediction, 'latest\_nontrivial\_error': final\_soc - latest\_prediction\} を返す。
            \end{enumerate}

\subsection*{L228 関数 \path|evaluate_soc_guard_intervention|}
            \begin{itemize}[leftmargin=1.6em]
              \item 定義: \path|evaluate_soc_guard_intervention(detail, manifest)|
              \item docstring: Reject policies that only finish because the execution safety guard intervened.
              \item このブロックが直接呼ぶ主な関数・メソッド: astype, bool, fillna, float, get, int, isin, lower, strip, sum, to\_numeric
              \item 戻り値の要点: \{'intervention\_rows': intervention\_rows, 'intervention\_sec': intervention\_sec, 'passed': bool(intervention\_rows == 0 and intervention\_sec <= 1e-09)\}
            \end{itemize}
            \begin{enumerate}[leftmargin=1.8em]
            \item 条件 'soc\_guard\_intervened' in detail.columns を判定し、真なら内部処理を行う。
\item \{'intervention\_rows': intervention\_rows, 'intervention\_sec': intervention\_sec, 'passed': bool(intervention\_rows == 0 and intervention\_sec <= 1e-09)\} を返す。
            \end{enumerate}

\subsection*{L249 関数 \path|main|}
            \begin{itemize}[leftmargin=1.6em]
              \item 定義: \path|main()|
              \item docstring: 明示されていない。
              \item このブロックが直接呼ぶ主な関数・メソッド: ArgumentParser, DataFrame, FileNotFoundError, Path, abs, add\_argument, all, append, bool, copy2, dumps, evaluate\_event\_timing
              \item 戻り値の要点: 0 if selection['selected'] else 2
            \end{itemize}
            \begin{enumerate}[leftmargin=1.8em]
            \item parser に argparse.ArgumentParser(description=\_\_doc\_\_) の結果を代入する。
\item parser.add\_argument(...) を実行する。
\item parser.add\_argument(...) を実行する。
\item parser.add\_argument(...) を実行する。
\item parser.add\_argument(...) を実行する。
\item parser.add\_argument(...) を実行する。
\item parser.add\_argument(...) を実行する。
\item parser.add\_argument(...) を実行する。
\item parser.add\_argument(...) を実行する。
\item args に parser.parse\_args() の結果を代入する。
            \end{enumerate}







        \subsection*{CLI 引数}
            \begin{longtable}{p{0.07\linewidth} p{0.78\linewidth}}
            \toprule
            行 & 引数 \\
            \midrule
            251 & \path|--profile| \\
252 & \path|--campaign-dir| \\
253 & \path|--output-dir| \\
254 & \path|--stage| \\
255 & \path|--seed-label| \\
260 & \path|--integration-ds-km| \\
261 & \path|--control-ds-km| \\
262 & \path|--no-resume| \\
            \bottomrule
            \end{longtable}



        \section*{処理の流れ}
        \begin{enumerate}[leftmargin=1.7em]
        \item CLI 引数を解釈し、main() から処理を起動する。
        \end{enumerate}

        \section*{このファイルの位置づけ}
        この資料群では、\path|scripts/validate_gpu_upper_policy_candidates.py| を
        \textbf{planning validation}の中核として扱う。
        したがって、ここで扱う処理は単独で完結せず、
        前段の profile / launch / telemetry と後段の planner / logger / report へ繋がる。

        \end{document}
